\clearpage
\section*{Appendix T. Necessary exterior-to-core energy transport}
\addcontentsline{toc}{section}{Appendix T. Necessary exterior-to-core energy transport}
\noindent\textbf{Scope.} This appendix sharpens the energy-compatibility
boundary of Appendix R while permitting a globally non-negligible
exterior correction. Its transport law is a necessary consequence of
the classical smooth unforced Navier--Stokes energy equation; neither
a positive H8 stress nor a singular unforced solution is constructed.

\paragraph{Local conservation and the moving boundary.}
Let $u$ be a smooth periodic divergence-free solution at unit viscosity,
and let $\pi$ be its physical periodic pressure. Put
$e=\frac12|u|^2$. The exact pointwise identity is
\begin{equation}\label{eq:T-local-energy}
 \partial_t e+
 \nabla\cdot\bigl[(e+\pi)u-\nabla e\bigr]
 =-|\nabla u|^2.
\end{equation}
Let $\Omega(t)\Subset\mathbb T^3$ have a smooth moving boundary,
outward normal $n$ and normal boundary velocity $V_b\cdot n$.
Reynolds transport gives
\begin{equation}\label{eq:T-moving-domain}
 \frac{d}{dt}\int_{\Omega(t)}e\,dx+
 \int_{\Omega(t)}|\nabla u|^2\,dx+
 \Phi(t)=0,\quad
 \Phi(t)=\int_{\partial\Omega(t)}
 \left[(e+\pi)u-\nabla e-eV_b\right]\cdot n\,dS.
\end{equation}
Here $\Phi>0$ denotes \emph{outward} transport. The
integrated pressure flux is invariant under the time-dependent
gauge $\pi\mapsto\pi+c(t)$ because
$\int_{\partial\Omega}u\cdot n\,dS=0$.

\paragraph{Theorem T.1 (Terminal-dyad inward-flux requirement).}
Let $\tau=T-t$, $p>0$ and $M_0,R_0>0$.
Assume a smooth prepared divergence-free field $U(t)$ is compactly
supported strictly inside $\Omega(t)$, vanishing in a collar of its
boundary, and satisfies
\begin{equation}\label{eq:T-U-scaling}
 \|U\|_2^2=M_0\tau^p(1+o(1)),\qquad
 \|\nabla U\|_2^2=R_0M_0\tau^{p-1}(1+o(1)).
\end{equation}
Suppose the genuine unforced solution $u=U+w$ retains the carrier
\emph{locally} in the moving core:
\begin{equation}\label{eq:T-local-L2}
 \|w(t)\|_{L^2(\Omega(t))}=o(\|U(t)\|_2).
\end{equation}
Assume additionally the dimensionless Laplacian shape ratio
$\|U\|_2\|\Delta U\|_2/\|\nabla U\|_2^2$ remains bounded.
Then for $I_\tau=[T-\tau,T-\tau/2]$,
\begin{equation}\label{eq:T-inward-lower}
 \boxed{\liminf_{\tau\downarrow0}
 \frac{-\int_{I_\tau}\Phi(s)\,ds}{M_0\tau^p}
 \ge \frac{1-2^{-p}}{p}\left(R_0-\frac p2\right).}
\end{equation}
If in addition
$\|\nabla w(t)\|_{L^2(\Omega(t))}
=o(\|\nabla U(t)\|_2)$,
then \eqref{eq:T-inward-lower} holds with
\emph{equality in the limit}.
Consequently $R_0>p/2$ requires a positive, leading-order
inward energy transfer across the moving boundary.
The theorem does not require $w$ to be small in global $L^2$.

\paragraph{Proof.}
Since $U$ is supported within $\Omega(t)$, the local retention
condition implies
\[
 E_\Omega(t):=\int_{\Omega(t)}e\,dx
 =\frac12 M_0\tau^p(1+o(1)).
\]
The collar permits integration by parts without surface terms:
\[
 \langle\nabla u,\nabla U\rangle_{L^2(\Omega)}
 =\|\nabla U\|_2^2+\langle w,-\Delta U\rangle_{L^2(\Omega)}.
\]
By Cauchy--Schwarz and the bounded shape ratio,
$\|\nabla u\|_{L^2(\Omega)}^2
\ge(1-o(1))\|\nabla U\|_2^2$.
The stronger gradient-retention condition makes this
asymptotic an equality.
Integrating \eqref{eq:T-moving-domain} over $I_\tau$ gives
\[
 -\int_{I_\tau}\Phi(s)\,ds
 =E_\Omega(T-\tau/2)-E_\Omega(T-\tau)
 +\int_{I_\tau}\!\int_{\Omega(s)}
 |\nabla u|^2\,dx\,ds.
\]
The first term is
$-\frac12M_0(1-2^{-p})\tau^p+o(\tau^p)$;
the second is at least
$(1-o(1))R_0M_0(1-2^{-p})\tau^p/p$.
Together they prove \eqref{eq:T-inward-lower}.
\hfill$\square$

\paragraph{Corollary T.2 (Source-independent axisymmetric diagnostic).}
For the compact axisymmetric solenoidal Appendix O prepared
profile, $h=1/10$, $p=1/5$, $X_*=1$, $|\eta|\le1/2$.
The transverse disk radius obeys $R_*^2=8/3$.
Because the horizontal Cartesian components of an axisymmetric
vector have angular order one and the axial component has zero
horizontal mean, the exact disk inequalities in Appendix O imply
\[
 R_0\ge\frac{3}{8}j_{1,1}^2
 =5.50573899079645\ldots .
\]
Therefore any unforced solution satisfying the local-retention
hypotheses of Theorem T.1 for this diagnostic must obey
\begin{equation}\label{eq:T-diagnostic-influx}
 \liminf_{\tau\downarrow0}
 \frac{-\int_{I_\tau}\Phi(s)\,ds}{M_0\tau^{1/5}}
 \ge \frac{1-2^{-1/5}}{1/5}
 \left(\frac38j_{1,1}^2-\frac1{10}\right)
 =3.49884933663390\ldots .
\end{equation}
The Bessel eigenvalue is classical; the numerical display is
a rounding of an exact spectral expression, not a certificate
for the non-authentic Appendix O source coefficients.

\paragraph{Proposition T.3 (An isolated compact core cannot retain).}
Under the hypotheses of Theorem T.1 with $R_0>p/2$,
if $w=0$ in an open collar of $\partial\Omega(t)$ for all
sufficiently late times, then no such unforced solution exists.
Indeed $U=0$ in that collar, hence $u=\nabla u=e=0$
on the moving boundary and \eqref{eq:T-moving-domain} has
$\Phi=0$, contradicting \eqref{eq:T-inward-lower}.
More generally, any proposed boundary estimates making
$\int_{I_\tau}|\Phi(s)|\,ds=o(\tau^p)$ contradict the same theorem.

\paragraph{Corollary T.4 (A bounded exterior collar cannot supply the flux).}
Assume also the H8 compact chart has lateral scale
$r=O(\tau^{1/2})$, axial scale $z=O(\tau^D)$,
where $D=1/2-h$, $0<h<1/6$, and $p=1/2-3h$.
Choose a smooth moving collar boundary with surface area
$|\partial\Omega(t)|=O(\tau^{1-h})$ and normal speed
$|V_b\cdot n|=O(\tau^{-A})$, $A=1/2+h$.
If the full unforced pressure $\pi$ and all traces
$w,\nabla w$ are uniformly bounded on this collar
as $t\uparrow T$, then, since $U=\nabla U=0$ there,
\[
 |\Phi(t)|=O(\tau^{1-h-A})
          =O(\tau^{1/2-2h}),\qquad
 \int_{I_\tau}|\Phi(s)|\,ds
          =O(\tau^{3/2-2h})
          =o(\tau^{p}).
\]
This contradicts Theorem T.1 whenever $R_0>p/2$.
Consequently, any exterior-coupled realization retaining such
a rapidly dissipating core must violate at least one of these
uniform boundary trace bounds: the correction or its gradient
or the physical pressure must develop a sufficiently large
interface contribution. This conditional exclusion does not
identify which channel supplies it, and does not require an
isolated or globally small correction.

\paragraph{A quantitative global-error alternative.}
Let $\epsilon(t)=\|u(t)-U(t)\|_2/\|U(t)\|_2$ globally,
$K=\limsup_{t\uparrow T}
 \|U\|_2\|\Delta U\|_2/\|\nabla U\|_2^2<\infty$,
and $\delta=\limsup_{t\uparrow T}\epsilon(t)$.
If $0\le\delta<\min\{1,1/K\}$, then the exact
global unforced energy law and the same projection argument give
the necessary inequality
\begin{equation}\label{eq:T-global-delta}
 \boxed{(1+\delta)^2-2^{-p}(1-\delta)^2
 \ge \frac{2R_0}{p}(1-2^{-p})(1-K\delta)^2.}
\end{equation}
To prove it, bound the early dyad energy above by
$(1+\delta+o(1))^2M_0\tau^p$ and the late energy
below by $(1-\delta+o(1))^2M_0 2^{-p}\tau^p$.
Meanwhile $\|\nabla u\|_2^2\ge
(1-K\delta-o(1))^2\|\nabla U\|_2^2$.
Insert these into the energy identity and pass to the limit.
At $\delta=0$ this recovers Theorem R.1's
necessary bound $R_0\le p/2$.
No numerical value for the authentic H8 ratio $K$ is asserted.

\paragraph{Where the exterior correction enters.}
Since $U$ and all its derivatives vanish on the selected
boundary collar, the integrand of the exact flux becomes
\[
 \left[\left(\frac{|w|^2}{2}+\pi\right)w
 -\nabla\left(\frac{|w|^2}{2}\right)
 -\frac{|w|^2}{2}V_b\right]\cdot n.
\]
Its nonlocal pressure trace is essential; neither the
amplitude and translation ODEs nor the pointwise stress cone
provides the required sign or size.
The decisive next source-specific certificate is a lower bound
on this \emph{actual} inward pressure--convection--viscous--boundary
flux on every sufficiently late dyad while independently proving
the complementary PDE and pointwise terminal retention.
The authentic H8 source tuple remains unfrozen.
The complete unforced Navier--Stokes theorem remains OPEN.
